\documentclass[10pt,a4paper]{article}
\usepackage[margin=1.6cm]{geometry}
\usepackage{amsmath, amssymb, graphicx, booktabs, paracol, xcolor}
\usepackage[hidelinks]{hyperref}
\setlength{\parindent}{0pt}
\setlength{\parskip}{0.45em}
\setlength{\columnsep}{1.2em}
\columnratio{0.5}
\graphicspath{{week_figure/fpca/}}
\newcommand{\fig}[2]{\begin{center}\includegraphics[width=\linewidth]{#1}\\[-0.2em]{\footnotesize #2}\end{center}}
\newcommand{\code}[1]{\texttt{\small #1}}
\newcommand{\inwords}[1]{\par\noindent\colorbox{gray!12}{\parbox{\dimexpr\linewidth-2\fboxsep\relax}{\textbf{In words.} #1}}\par}
\newcommand{\maths}[1]{\par\noindent\colorbox{blue!5}{\parbox{\dimexpr\linewidth-2\fboxsep\relax}{\footnotesize\textbf{Formula.} #1}}\par}

\title{FPCA: the main ways people differ\\[0.2em]\large How the curvature shapes are found, idea by idea}
\date{4 November 2026}

\begin{document}
\sloppy
\maketitle
\vspace{-1em}
\begin{paracol}{2}

Code: \code{tortuosity/curvature.py} (\code{profile\_thesis}) and
\code{CGPS\_data/turning\_profile\_shapes.py} (\code{pca}, \code{shape\_count}). Figures use 300
made-up people run through the same code, except the last (CGPS counts). Grey boxes give the idea,
blue boxes the formula; the text can be read without the blue boxes.

\switchcolumn
\switchcolumn*

%=======================================================================================
\section{What goes in: one curve per person}
\inwords{At every point of the segment, the turning is how sharply the vessel changes direction over
5\,mm. Each segment is stretched or squeezed to run from 0 (start) to 1 (end) and read off at the same
50 places, so place 15 means ``30\,\% of the way along'' for everyone. Each person is then one curve of
50 numbers, and all people together are a table $M$ with one row per person.}

Bends shorter than about 5\,mm are not seen; the 5\,mm is the scale. The first and last 5\,mm of a
trunk get no value. There is one table per segment.

\maths{$\theta$ is the angle between the 5\,mm chord arriving at a point and the one leaving it;
$T = \theta/5\,\text{mm}$ (rad/mm). Places $u_j = (j - 0.5)/50$, linear interpolation.}

\switchcolumn
\fig{f1_turning_geometry}{Figure 1. The turning at one point: the angle between the chord arriving and
the chord leaving.}
\fig{f3_grid}{Figure 2. One segment in mm (top) and on the 50 places (bottom).}
\switchcolumn*

%=======================================================================================
\section{Remove what everybody shares}
\inwords{The average curve says how much everybody turns at each place, for example that everybody
turns more near the start. It is the same for everyone, so it cannot tell people apart. At each place
the average there is subtracted from everyone. What is left, the deviation, is how much and where each
person turns more or less than the average person.}

Subtracting the same number from everyone at a place keeps their order and the gaps between them
(Figure~3, bottom). A deviation above 0 means the person turns more than average at that place, below 0
less.

\textbf{At each place the deviations add up to 0}, as those above the average cancel those below. Two
things follow. A score of 0 below means ``like the average person''. And without the subtraction the
first pattern found would be the average shape itself (``everybody turns like this''), not a way people
differ.

\maths{$m_j = \tfrac{1}{n}\sum_i M_{ij}$, \; $D_{ij} = M_{ij} - m_j$, \; $\sum_i D_{ij} = 0$. In the code:
\code{mean = M.mean(0)}, \code{M - mean}.}

\switchcolumn
\fig{f4_profiles_deviations}{Figure 3. Top: curves of 25 of the 300 people and the average (black).
Bottom: deviations; the width of the bundle at a place is how much people differ there.}
\switchcolumn*

%=======================================================================================
\section{A pattern and a score}
\inwords{A pattern is a curve of weights along the segment, one weight per place. A person's score is
how much of the pattern they have. It is found by multiplying the person's deviation by the pattern's
weight at every place and adding the 50 products up. Where both have the same sign the score rises,
where they have opposite signs it falls, and where the weight is near 0 the place hardly counts.}

The person in Figure~4 turns much less than average near the start and about average further on.
\begin{itemize}\setlength\itemsep{0pt}
  \item $\varphi_1$ has positive weight everywhere (``turns more everywhere''). The person's deviation is
        mostly negative, so the products are mostly negative: score $-1.6$\,SD, less turning overall.
  \item $\varphi_2$ is positive at the start and negative at the end (``more at the start, less at the
        end''). The large negative deviation at the start gives a large negative product: score
        $-2.3$\,SD, the opposite of the pattern, i.e.\ less at the start relative to the end.
\end{itemize}

A large positive score means the person clearly has the pattern, a large negative score its opposite,
near 0 that they do not differ from the average in that way. Every pattern has the same size (the
squared weights add to 1), so scores along different patterns are on the same footing.

The score is also the amount of the pattern that, added to the average curve, comes closest to the
person's curve (worked through in \code{shape\_score\_explained.tex}).

\maths{$c_{ik} = \sum_j D_{ij}\varphi_{kj}$, \; $\sum_j \varphi_{kj}^2 = 1$.}

\switchcolumn
\fig{f6b_score}{Figure 4. Top: one person's deviation. Middle: the first two patterns. Bottom: deviation
times weight at every place; the sum of the coloured minus the grey area is the score (shown in SD of all
people's scores).}
\switchcolumn*

%=======================================================================================
\section{The first pattern: the one that spreads people most}
\inwords{Any curve of weights gives every person a score. A useful pattern is one that tells people
apart: their scores along it are far apart. The first pattern, $\varphi_1$, is the curve of weights
along which people's scores spread most. No plaque information is used.}

\textbf{Spreading most = missing least.} Each person's deviation splits into two parts: the part the
pattern describes (score times pattern) and the part it misses. The two parts add up, in squares, to the
person's whole deviation, which is fixed by the data. A pattern that gives large scores therefore leaves
little unexplained. $\varphi_1$ is the single curve that, on its own, describes most of how people
differ.

The spread along $\varphi_1$ has a name, the eigenvalue $\lambda_1$: the variance of the scores along it.

In the made-up data (Figure~5), scores along $\varphi_1$ spread most (SD 0.042), along $\varphi_2$ less,
along $\varphi_3$ least.

\maths{$\lambda_k = \tfrac{1}{n}\sum_i c_{ik}^2$; \;
$\sum_j D_{ij}^2 = c_{ik}^2 + \sum_j (D_{ij} - c_{ik}\varphi_{kj})^2$ (Pythagoras).}

\switchcolumn
\fig{f6c_spread}{Figure 5. Scores of all 300 people along the first three patterns. Each pattern spreads
people less than the one before.}
\switchcolumn*

%=======================================================================================
\section{Remove it and ask again}
\inwords{Subtract from each person the part $\varphi_1$ describes: their score times the pattern. What is
left has nothing of $\varphi_1$ in it any more. Ask the same question on what is left: which curve of
weights spreads people most now? That is $\varphi_2$. Repeat.}

Figure~6 shows the three rounds. In round 1, $\varphi_1$ (``more everywhere'') is found and removed; 42\,\%
of the differences are left. In round 2, $\varphi_2$ (``start against end'') is found in what is left,
and so on.

\textbf{Each new pattern shares nothing with the earlier ones.} Nothing of $\varphi_1$ is left, so
$\varphi_2$ cannot contain any of it: multiplying their weights place by place and adding gives 0. The
patterns therefore never describe the same thing twice, and each person's deviation splits into separate
pieces, one per pattern.

There are as many patterns as places, 50. Most of the later ones describe noise (Section~\ref{sec:k}).

\maths{left after round 1: $D_{ij} - c_{i1}\varphi_{1j}$; \; $\sum_j \varphi_{1j}\varphi_{2j} = 0$.}

\switchcolumn
\fig{f6_steps}{Figure 6. Find a pattern, remove it from everyone, repeat. Left: what is not yet
described and the pattern found ($\pm 2$\,SD). Right: what is left after removing it.}
\switchcolumn*

%=======================================================================================
\section{How much each pattern describes}
\inwords{The spreads of all patterns add up to the total difference between people. A pattern's share is
its spread divided by the total. A person can be rebuilt as the average plus their score times each
pattern; each added pattern brings the rebuilt curve closer to the person.}

In the made-up data $\varphi_1$ describes 58\,\%, $\varphi_2$ 29\,\% and $\varphi_3$ 12\,\% of how people
differ. These shares are what the scree plot shows (Figure~10).

Figure~7 rebuilds the person of Figure~4. The average alone misses all of their deviation; adding
$\varphi_1$ lowers the whole curve; adding $\varphi_2$ lowers the start further; with $\varphi_3$ almost
nothing is left. The made-up data were built from three patterns plus noise, so three patterns are
enough here.

\maths{share$_k = \lambda_k / \sum_{k'} \lambda_{k'}$; \; rebuilt curve $= m + c_1\varphi_1 + c_2\varphi_2
+ \dots$}

\switchcolumn
\fig{f8a_rebuild}{Figure 7. The person of Figure~4 (black) rebuilt from the average plus 0, 1, 2 and 3
patterns. Grey: the part not yet described.}
\switchcolumn*

%=======================================================================================
\section{Why the patterns are smooth, and how they are computed}
\inwords{A bend covers several neighbouring places, so a person who turns more than average at one place
also turns more at its neighbours. The patterns therefore come out as smooth curves, not as jumps from
place to place. The computer does not try curves one by one; it finds all patterns at once from how the
places move together.}

Figure~8 (left) shows how places move together: red where people who turn more than average at one
place also turn more at the other. It is red near the diagonal, where places are neighbours. The
patterns (right) are read off from this table.

FPCA finds the ways people \emph{differ} most, not the ways that matter for plaque.

\maths{$S_{jj'} = \tfrac{1}{n}\sum_i D_{ij}D_{ij'}$. The patterns are the eigenvectors of $S$
($S\varphi = \lambda\varphi$), the spreads its eigenvalues. The code gets the same patterns from a
singular value decomposition of $D$ (\code{np.linalg.svd(M - mean)}).}

\switchcolumn
\fig{f5_covariance}{Figure 8. Left: how much the deviations at two places move together. Right: the first
three patterns found from it, with their shares.}
\switchcolumn*

%=======================================================================================
\section{Reading a pattern}
\inwords{A pattern is easiest to read as two people: one with a lot of it and one with its opposite,
drawn around the average. Where the weights are positive the first turns more than average; where the
pattern crosses 0 the two swap sides. Each pattern is named in words before plaque is looked at.}

In Figure~9, $\varphi_1$ is ``more everywhere'', $\varphi_2$ ``more at the start, less at the end'' and
$\varphi_3$ ``more in the middle, less at the ends''.

\textbf{The sign is a free choice.} A pattern and its opposite describe the same difference; only the
sign of the scores flips. The code fixes it so that each pattern's weights add up to $\ge 0$.

\textbf{Standardised score} $z = c/\mathrm{SD}$: $+1$ means one SD more of the pattern than the
average person. This is what enters the model.

\maths{high: $m + 2\,\mathrm{SD}_k\,\varphi_k$; low: $m - 2\,\mathrm{SD}_k\,\varphi_k$.}

\switchcolumn
\fig{f7_mean_pm}{Figure 9. Each pattern as a person with a lot of it (solid) and with its opposite
(dashed), around the average.}
\switchcolumn*

%=======================================================================================
\section{What comes out}
\label{sec:k}
\textbf{How many patterns are real.} Three checks, none of which looks at plaque; $K$ is fixed before
the outcome is linked.
\begin{itemize}\setlength\itemsep{0pt}
  \item Horn: shuffle each place over people, so no real pattern is left, and redo the FPCA 200 times.
        A pattern is real if its share beats the 95th percentile of the shuffled shares.
  \item Bootstrap: redraw the people with replacement 200 times and redo the FPCA. A pattern is stable
        if a redrawn pattern matches it closely: median best $|\text{cosine}| \ge 0.9$.
  \item Cross-validation: hide one fifth of a held-out person's curve (one block of places), estimate
        the scores from the rest and predict the hidden part. $K$ is where the error is lowest.
\end{itemize}
$K$: the patterns that pass Horn and the bootstrap, capped by the budget below.

\switchcolumn
\fig{f9_scree}{Figure 10. Share per pattern against chance. Made-up data built from three patterns: three
pass.}
\fig{f10_boot_cv}{Figure 11. Bootstrap stability and cross-validation error.}
\switchcolumn*

\textbf{The descriptive pLAD and mLAD models.} The main who model uses the mean turning per segment.
The shapes are a descriptive analysis in the proximal and mid LAD only, one model per segment, not
pooled and not part of the prespecified test. Each model has SCORE2 and the standardised scores
$z_1, \dots, z_K$ and is not coded yet:
\[
  \log\!\bigl(-\log(1 - p_i)\bigr) = a + g\,\mathrm{SCORE2}_i + b_1 z_{i1} + \dots + b_K z_{iK},
\]
$p_i$ the risk that person $i$ develops plaque in the segment. $e^{b_k}$ is the hazard ratio per SD of
pattern $k$. With $z_1$ in the model, $b_2, \dots, b_K$ ask whether the \emph{position} of the bends
matters beyond how much the segment turns. The weights combine into one weight per place,
$\beta_j = \sum_k (b_k/\mathrm{SD}_k)\,\varphi_{kj}$.

\textbf{Budget.} With 275 (pLAD) and 201 (mLAD) people who developed plaque in the segment, Riley's
criteria allow $K = 5$ and 4 at $C = 0.65$; the other segments have room for at most 2.

\switchcolumn
\fig{f11_model}{Figure 12. Top: how the model's sum becomes a risk. Bottom: the weight per place for
example weights $b = (0.15, -0.10, 0)$.}
\fig{f12_budget}{Figure 13. $K_{\max}$ per segment (CGPS, 1,077 plaque-free participants, before the
image-quality exclusion).}
\switchcolumn*

\end{paracol}
\end{document}
